\documentclass[11pt,a4paper]{article}

% === Packages =============================================
\usepackage[utf8]{inputenc}
\usepackage[T1]{fontenc}
\usepackage{lmodern}
\usepackage{amsmath,amssymb}
\usepackage{graphicx}
\usepackage{booktabs}
\usepackage{hyperref}
\usepackage{xcolor}
\usepackage{geometry}
\usepackage{caption}
\usepackage{subcaption}
\usepackage{cite}
\usepackage{microtype}

\geometry{margin=2.5cm}

\hypersetup{
  colorlinks=true,
  linkcolor=blue,
  citecolor=blue,
  urlcolor=blue
}

% === Title ===============================================
\title{
  \textbf{Semantic Bundle AI: Anchor Design and Stability Framework} \\[0.5em]
  \large Stable Semantic Regions, Domain Discriminability, \\
  and a Four-Axis Evaluation Framework for Semantic Stability
}

\author{
  M. Saitou \\
  \textit{ON THE SHOULDERS OF GIANTS PROJECT} \\
  \texttt{m.saitou@ontheshouldersofgiants.jp}
}

\date{May 2026\thanks{This work is licensed under CC BY-NC-ND 4.0.
\url{https://creativecommons.org/licenses/by-nc-nd/4.0/}}}

% === Body ================================================
\begin{document}

\maketitle

\begin{abstract}
Semantic Bundle AI introduces stable anchor-based coordinates as a
complementary semantic management layer for Large Language Models
\cite{saitou2026semantic, saitou2026poc}. A key open question from
prior work concerns anchor design: what constitutes an effective anchor,
and how should anchors be selected?

This paper addresses that question through two experiments. First, we
demonstrate that anchors constructed as \textit{stable semantic regions}
--- averaged embeddings of meaning-neighbor clusters rather than single
words --- achieve substantially higher inter-model stability than
word-level anchors (Spearman $r = 0.704$ vs.\ $0.101$, a 7$\times$
improvement). Second, we identify a fundamental tradeoff between cluster
stability and domain discriminability: generic anchors excel at
intra-cluster cohesion (variance ratio $0.047$--$0.067$), while
domain-specific concept vectors achieve $4.4\times$ higher inter-domain
separation ($0.739$ vs.\ $0.166$). We propose a four-axis evaluation
framework for semantic stability --- inter-model consistency, temporal
consistency, cross-cultural robustness, and perturbation resistance ---
and provide initial validation on two axes. These results establish
concrete design principles for anchor selection and position semantic
stability as an empirically tractable research program.
\end{abstract}

% --------------------------------------------------------
\section{Introduction}
\label{sec:intro}

Prior work on Semantic Bundle AI established the theoretical framework
\cite{saitou2026semantic} and provided initial empirical validation across
four axes: anchor coordinate stability, longitudinal consistency, edit
locality, and compression efficiency \cite{saitou2026poc}. A central
finding was that anchor coordinates reduce intra-cluster variance to 5.2\%
of raw embedding variance, while bundle updates reduce semantic drift by
38.6\% over ten sequential updates.

However, a critical question was left open: \textit{how should anchors be
designed?} The prior work used generic concept words (time, space, risk,
rule, trust, number, body, emotion) without systematic justification. The
failure of hypothesis H-102 --- cross-model ranking consistency did not
improve under anchor projection --- suggested that anchor design may be
the limiting factor.

This paper addresses anchor design through two lines of investigation:

\begin{enumerate}
\item \textbf{Anchor representation}: Should anchors be single words or
richer semantic structures? We compare word-level anchors, stable semantic
region anchors (meaning-neighbor clusters), and embodied cognition anchors
(body-based conceptual primitives).

\item \textbf{Stability--discriminability tradeoff}: Is there a fundamental
tradeoff between using anchors for intra-cluster cohesion versus
inter-domain separation? We characterize this tradeoff and derive
application-specific design recommendations.
\end{enumerate}

We further propose a four-axis framework for semantic stability evaluation,
providing a principled structure for future research.

% --------------------------------------------------------
\section{Background}
\label{sec:background}

\subsection{Anchor Design in Semantic Spaces}

The concept of using reference points to stabilize semantic representations
has appeared in several contexts. Procrustes alignment \cite{smith2017offline}
uses reference words to align embedding spaces across languages. Diachronic
embedding studies \cite{hamilton2016diachronic} use stable landmark words
to track semantic change over time. These approaches treat anchors as
single words selected by frequency or stability criteria.

Semantic Bundle AI differs in motivation: anchors are not used for
cross-lingual alignment or drift measurement, but as an ongoing coordinate
reference for semantic management. This distinction has implications for
anchor design that have not been previously explored.

\subsection{Embodied Cognition and Semantic Stability}

A theoretical motivation for body-based anchor selection comes from
embodied cognition research \cite{lakoff1999philosophy}, which argues
that abstract concepts are grounded in sensorimotor experience. Concepts
such as temperature (hot/cold), proximity (near/far), and verticality
(up/down) are hypothesized to be more cross-culturally stable because
they are anchored in shared human embodiment. This motivates testing
whether embodied concept clusters yield more stable anchors.

\subsection{Semantic Region vs.\ Word-Level Representation}

Rather than treating anchors as individual word vectors, we propose
constructing anchors as \textit{stable semantic regions}: averaged
embeddings of meaning-neighbor clusters that include synonyms,
multilingual equivalents, and contextually related terms. This is
motivated by the observation that single-word embeddings are sensitive
to tokenization, frequency, and context, while averaged representations
over semantic neighborhoods tend to be more stable \cite{wieting2016towards}.

% --------------------------------------------------------
\section{Experimental Setup}
\label{sec:setup}

\subsection{Models}

All experiments use two pre-trained sentence embedding models without
modification: \texttt{all-mpnet-base-v2} ($d = 768$) and
\texttt{all-MiniLM-L6-v2} ($d = 384$) from the
\texttt{sentence-transformers} library \cite{reimers2019sentence}.

\subsection{Anchor Types}

We compare three anchor construction strategies, each with 8 anchors:

\textbf{Word anchors}: Single-word embeddings for generic concepts
(time, space, risk, rule, trust, number, body, emotion).

\textbf{Semantic region anchors}: Each anchor is the mean embedding of
10 semantically related terms spanning synonyms and multilingual
equivalents. For example, the \textit{danger} anchor averages embeddings
of \{danger, risk, unsafe, hazardous, threat, \textit{kiken}, \textit{risuku}, harm, peril,
jeopardy\}.

\textbf{Embodied anchors}: Each anchor is the mean embedding of 10
body-based conceptual primitives. Eight embodied concepts are used:
hot/cold, near/far, up/down, heavy/light, safe/danger, fast/slow,
many/few, past/future.

\subsection{Evaluation Metrics}

\textbf{A-axis (inter-model stability)}: Spearman rank correlation
between the pairwise similarity matrices of anchor embeddings computed
by the two models. High values indicate that the relational structure
among anchors is preserved across models.

\textbf{Cluster stability} (H-101 equivalent): Ratio of intra-cluster
variance in anchor coordinate space to raw embedding space for a
synonym group of 15 sentences. Lower values indicate tighter clustering.

\textbf{D-axis (perturbation resistance)}: Correlation between full and
anchor-reduced coordinate representations when one anchor is removed.
Values near 1.0 indicate robustness to anchor removal.

\textbf{H-603 (ranking consistency)}: Spearman rank correlation of
nearest-neighbor rankings between the two models in anchor coordinate
space, following the H-102 protocol from \cite{saitou2026poc}.

\subsection{Domain Discriminability Experiment}

We construct four domain datasets of 20 sentences each: information
security, medical, legal, and general. Two anchor types are compared
for domain discriminability: generic word anchors and domain-specific
concept vectors. Inter-domain separation is measured as the mean
Euclidean distance between domain centroids in coordinate space.

% --------------------------------------------------------
\section{Results}
\label{sec:results}

\subsection{Anchor Type Comparison}

Table~\ref{tab:anchor_comparison} summarizes results across all metrics
for the three anchor types.

\begin{table}[h]
\centering
\caption{Comparison of anchor types across evaluation metrics.}
\label{tab:anchor_comparison}
\begin{tabular}{lrrrr}
\toprule
Anchor type & A-axis & Cluster ratio & D-axis & H-603 \\
\midrule
Word (single)         & 0.101 & \textbf{0.036} & 1.000 & 0.617 \\
Semantic region       & \textbf{0.704} & 0.047 & 1.000 & 0.643$^\dagger$ \\
Embodied (body-based) & 0.559 & 0.042 & 1.000 & 0.596 \\
\midrule
Paper 1 baseline      & ---   & 0.052 & ---   & 0.572--0.643 \\
\bottomrule
\end{tabular}

\smallskip
\noindent{\footnotesize $^\dagger$ H-603 for semantic region matches the Paper~1 baseline upper bound (no improvement).}
\end{table}

Several findings emerge. First, semantic region anchors achieve
substantially higher inter-model stability (A-axis: $0.704$) compared
to word anchors ($0.101$), a $7\times$ improvement. Embodied anchors
occupy an intermediate position ($0.559$). Second, cluster stability
(variance ratio) is comparable across all three types ($0.036$--$0.047$),
suggesting that cluster cohesion is relatively insensitive to anchor
construction strategy. Third, perturbation resistance (D-axis) is
perfect ($1.000$) for all types, replicating the H-103/H-104 findings
from \cite{saitou2026poc}. Fourth, H-603 ranking consistency for semantic region anchors
($0.643$) matches, without exceeding, the upper bound of the Paper 1 baseline range,
while word anchors fall below it ($0.617$); anchor projection thus does not improve
ranking consistency over the baseline.

Results are shown in Figure~\ref{fig:anchor_comparison}.

\begin{figure}[h]
  \centering
  \includegraphics[width=\textwidth]{figures/paper2_02_results_final.png}
  \caption{Anchor type comparison. \textit{Left}: Cross-model ranking
  consistency (H-603). \textit{Right}: All metrics comparison. Semantic
  region anchors (blue) lead on inter-model stability; among the three
  anchor types they also rank highest on ranking consistency, which
  nonetheless remains at the Paper~1 baseline level, while all types
  achieve comparable cluster stability and perfect perturbation resistance.}
  \label{fig:anchor_comparison}
\end{figure}

\subsection{Stability--Discriminability Tradeoff}

Table~\ref{tab:tradeoff} and Figure~\ref{fig:tradeoff} show results for
the stability--discriminability comparison across domains.

\begin{table}[h]
\centering
\caption{Stability vs.\ discriminability for generic vs.\ domain-specific anchors.}
\label{tab:tradeoff}
\begin{tabular}{lrr}
\toprule
Anchor type & Cluster stability (ratio) & Inter-domain separation \\
\midrule
Generic (word)        & \textbf{0.047--0.067} & 0.166 \\
Domain-specific       & 0.110--0.235          & \textbf{0.739} \\
\midrule
Improvement           & $\times$0.3--0.5 worse & $\times$4.4 better \\
\bottomrule
\end{tabular}
\end{table}

Generic anchors achieve lower cluster variance ratios ($0.047$--$0.067$)
than domain-specific concept vectors ($0.110$--$0.235$), indicating
better intra-cluster cohesion. However, domain-specific vectors achieve
$4.4\times$ higher inter-domain separation ($0.739$ vs.\ $0.166$).
This tradeoff is consistent across all four domains tested.

\begin{figure}[h]
  \centering
  \includegraphics[width=\textwidth]{figures/paper2_01_results.png}
  \caption{Stability--discriminability tradeoff. \textit{Left}: Cluster
  stability ratios by domain and anchor type. \textit{Center}: Inter-domain
  separation scores. \textit{Right}: Tradeoff scatter plot showing generic
  anchors (blue) favoring stability and domain-specific vectors (orange)
  favoring discriminability.}
  \label{fig:tradeoff}
\end{figure}

% --------------------------------------------------------
\section{A Four-Axis Framework for Semantic Stability}
\label{sec:framework}

Based on the experimental findings, we propose organizing semantic
stability evaluation around four axes:

\begin{equation}
\text{stability\_score}(a_i) =
  w_1 \cdot \text{inter\_model}(a_i) +
  w_2 \cdot \text{temporal}(a_i) +
  w_3 \cdot \text{cross\_cultural}(a_i) +
  w_4 \cdot \text{perturbation}(a_i)
\label{eq:stability}
\end{equation}

\textbf{A. Inter-model consistency}: Preservation of relational structure
across embedding models (validated in this work: $r = 0.704$ for semantic
region anchors).

\textbf{B. Temporal consistency}: Preservation of anchor position under
model updates and fine-tuning (partially validated in \cite{saitou2026poc}
via longitudinal bundle updates).

\textbf{C. Cross-cultural robustness}: Preservation of relational structure
across languages and cultural contexts (left for future work).

\textbf{D. Perturbation resistance}: Robustness to anchor removal and noise
injection (validated: $1.000$ across all anchor types).

The key contribution of this framework is to transform semantic stability
from an informal notion into an empirically tractable, multi-dimensional
evaluation problem. Axes B and C remain open research questions, providing
a structured agenda for future work.

% --------------------------------------------------------
\section{Discussion}
\label{sec:discussion}

\subsection{Design Principles for Anchor Selection}

The results support two practical design principles:

\textbf{Principle 1: Use semantic regions, not single words.}
Averaging over meaning-neighbor clusters yields $7\times$ higher
inter-model stability without sacrificing cluster cohesion or
perturbation resistance.

\textbf{Principle 2: Match anchor type to application.}
For applications requiring long-term semantic consistency (memory
management, concept tracking), generic anchors with lower cluster
variance ratios are preferable. For applications requiring domain
discrimination (classification, retrieval), domain-specific concept
vectors provide substantially higher inter-domain separation.

\subsection{Relationship to Embodied Cognition}

Embodied anchors occupy an intermediate position on all metrics,
suggesting partial support for the embodied cognition motivation.
They outperform word anchors substantially on inter-model stability
($0.559$ vs.\ $0.101$) but do not reach semantic region performance
($0.704$). This may reflect that body-based concepts, while cognitively
grounded, do not form as tight semantic neighborhoods in embedding
space as explicitly constructed meaning clusters. Further investigation
with larger embodied concept inventories is warranted.

\subsection{Implications for the H-102 Failure}

The H-102 failure in \cite{saitou2026poc} --- cross-model ranking
consistency did not improve under anchor projection --- can now be
partially explained. Word-level anchors impose an unstable inter-model
coordinate system ($r = 0.101$), whereas semantic region anchors are
far more stable ($r = 0.704$). This greater anchor stability is
associated with a modest gain in ranking consistency across anchor types
($0.617 \to 0.643$). However, even the most stable anchors only match
the upper bound of the Paper~1 baseline range ($0.643$) without
exceeding it; anchor design alone therefore does not recover cross-model
ranking consistency. The H-102 limitation persists and is consistent
with the cluster-stability/ranking tradeoff imposed by low-dimensional
anchor projection, rather than being resolved by anchor design.

\subsection{Limitations}

The current experiments use small datasets (15--20 sentences per domain)
and two embedding models. Axes B (temporal) and C (cross-cultural) of
the stability framework remain unvalidated. The weight parameters $w_i$
in Equation~\ref{eq:stability} are not yet determined empirically.
Large-scale validation across more models, languages, and domains is
required before deployment recommendations can be made with confidence.

% --------------------------------------------------------
\section{Conclusion}
\label{sec:conclusion}

We have presented two experiments addressing anchor design for Semantic
Bundle AI. Semantic region anchors --- constructed as averaged embeddings
of meaning-neighbor clusters --- achieve $7\times$ higher inter-model
stability than word-level anchors while maintaining comparable cluster
cohesion and perfect perturbation resistance. A fundamental
stability--discriminability tradeoff is identified: generic anchors
favor intra-cluster cohesion, while domain-specific concept vectors
achieve $4.4\times$ higher inter-domain separation.

We propose a four-axis framework for semantic stability evaluation,
providing initial validation on two axes and a structured agenda for
future work. Together, these results establish anchor design as a
tractable subproblem within the broader Semantic Bundle AI research
program, and provide concrete principles for practical anchor selection.

% --------------------------------------------------------
\bibliographystyle{plain}
\bibliography{references}

\end{document}
